\section{Mobility and transport impacts}
\label{sec:mobility_and_transport_impacts}

This Section reviews the main mobility impacts of \acp{bss}, focusing on how they shape trip patterns, interact with \ac{pt}, alter modal share, and influence broader transport system performance.

    \subsection{Trip patterns and multimodal integration}

Bike-sharing has measurably altered urban mobility patterns by adding a flexible, short-trip transport option that often complements existing modes. Evidence consistently shows that \acp{bss} are used predominantly for utilitarian purposes, with commuting and errands being the most common, while also supporting leisure and tourism mobility~\cite{Ricci2015}.

% Time efficiency gains
A frequently cited benefit of this utilitarian usage is improved first- and last-mile connectivity to \ac{pt}~\cite{Shaheen2010, Qiu2018}. By strategically siting bike stations near \ac{pt} hubs, \ac{bss} can extend the reach of metro and bus networks, and therefore facilitate multimodal trips~\cite{Cheng2022}. In Beijing, for example, bike-sharing access to bus and metro stations saved commuters an average of eight minutes per day, translating into substantial citywide time savings and even modest productivity gains~\cite{Qiu2018}. Recent analyses also point to additional efficiency gains from \acp{e-b}, whose travel speeds are approximately 9\% higher than those of \acp{m-b}, which is equivalent to roughly 70 seconds saved on a typical 1.5–2 km trip~\cite{Waldner2025}, further highlighting the contribution of contemporary \acp{bss} to urban time efficiency. 

    \subsection{Modal substitution and car-use reduction}

% Modal substitution
Beyond connectivity and time efficiency, \acp{bss} capacity for modal substitution is central, as it is one of the most debated aspects of their impact. Understanding which modes bike-sharing replaces is essential, as reducing car dependency is a core motivation for their introduction~\cite{Fishman2016}. Early surveys of the 2010s across cities including Dublin, London, Lyon, Barcelona, Montreal, Washington D.C., Minneapolis, Melbourne and Brisbane consistently showed that only 5--20\% of bike-sharing trips would otherwise have been made by car, with the majority substituting for walking or \ac{pt}~\cite{Murphy2015, Fishman2013, Fishman2014, BachandMarleau2012}. Ricci~\cite{Ricci2015} concluded that there was “no evidence that bike-sharing significantly reduces traffic congestion, carbon emissions and pollution”. At the same time, she cautioned that examining mode substitution alone does not capture the full range of impacts. To obtain a more complete assessment, she argued that research should also consider the characteristics of replaced car trips (e.g., frequency, duration, route) and the extent of induced demand through new trips that would not otherwise have been made—factors that were rarely examined in the earlier survey-based literature. On this basis, both Ricci and Fishman called for more detailed research on this issue in 2015~\cite{Ricci2015, Fishman2016}.  

Subsequent research broadly supports the finding that only a minority of bike-sharing trips replace car travel. Despite the technological advances of \ac{e-b} and dockless systems, the consensus remains the same: while bike-sharing does replace some car usage, the majority of trips continue to come from walking and \ac{pt}. In Beijing a survey of dockless system users found that only 14.8\% of trips replaced car travel~\cite{Chen2022b}; in Delft, the introduction of both docked (\textit{OV-fiets}) and dockless (\textit{Mobike}) systems showed similar patterns, with approximately 30\% of trips replacing car use~\cite{Ma2020}; and a multi-city European survey in 2021--2022 placed the figure at about 21\%~\cite{Bosehans2024}. Although these figures are higher than early estimates, they still indicate that most bike-sharing usage comes at the expense of active modes or \ac{pt} rather than driving, and that some authors argue that \ac{e-b} are not enough to attract car users to \acp{bss}~\cite{Teixeira2023}. Accordingly, some authors argue that more restrictive measures, such as car-free zones or increased parking charges, are necessary to strengthen the car-reducing effects of \acp{bss}~\cite{Bosehans2024}.

Even when bike-sharing does not fully replace car journeys, it can reduce car dependence indirectly by enabling multimodal travel. Many users combine bike-sharing with \ac{pt} for first- and last-mile segments, and repeated participation in bike-sharing has been associated with reductions in driving frequency. In Washington D.C., for example, 55\% of respondents to the \textit{CaBi} 2016 annual survey reported driving less often since joining the system~\cite{NACTO2020}.

    \subsection{Contextual factors shaping modal substitution}

% Modal substitution context matters
Context, however, matters. Research highlights that \ac{bss} user demographics and trip purposes shape modal substitution outcomes. Dockless systems in China and Europe often attract younger riders and those without cars or driving licenses, meaning these users were unlikely to be replacing car trips in the first place~\cite{Ma2020, Bosehans2024}. By contrast, in cities or niches where bike-sharing directly competes with car convenience, higher car substitution rates are possible. Moreover, cargo bike-sharing services--which allow users to haul goods or children--show particularly strong potential to displace car use: one survey of cargo bike-sharing users found that 46\% would have driven a car if the cargo bike had not been available~\cite{Becker2018}. This suggests that tailoring bike-sharing offerings (e.g. \acp{e-b} for longer commutes, cargo bikes for errands) can increase the likelihood of replacing car trips.

    \subsection{System-level effects of modal substitution}

% Congestion, parking and car ownership reduction
Although direct car substitution rates remain modest, even small shifts can yield measurable system-level benefits. In Washington D.C., \textit{Capital Bike-sharing} reduced neighborhood traffic congestion by about 4\%, with the greatest reductions in already congested areas~\cite{Hamilton2018}. A study of 96 U.S. cities similarly found that bike-sharing deployment eased congestion in large, transit-rich cities, particularly at rush hours, though effects were weaker or mixed in smaller and high-car-ownership contexts~\cite{Wang2017}. In China’s megacities such as Beijing, Shanghai, and Wuhan, the initial rollout of dockless \acp{bss} significantly reduced congestion in the short term, though subsequent expansions had diminishing returns~\cite{Xu2024}. Notably, in these Chinese cities, metro ridership rose alongside bike-sharing use, suggesting that many users shifted from driving to bike+metro combinations. Parking impacts follow a similar logic: in Pittsburgh, the opening of a \ac{bss} station was associated with a 2\% decline in parking meter use~\cite{Pelechrinis2016}, while in Boston the installation of new stations led to a 2.2\% reduction in household vehicle ownership and a 3.3\% drop in vehicle-miles traveled, with the strongest impacts (around 10\% lower car dependence) observed near transit stations~\cite{Basu2021}. Together, these findings suggest that \acp{bss}, as part of a multimodal system, can support “car-lite” lifestyles and free up curb space over the long term.  

% Induced demand
Beyond substitution, \acp{bss} also induce new trips that would not otherwise occur. A European study found that about 11\% of shared mobility trips were entirely new journeys induced by the service~\cite{Bosehans2024}. Such trips do not reduce any particular mode, but they indicate that bike-sharing provides additional mobility and access, expanding travel opportunities for its users. 

    \subsection{Transport network resilience}

% Transport network resilience
Finally, another key mobility impact of \acp{bss} is their contribution to the transport network resilience. During four London metro strikes, bike-sharing trips spiked sharply, with many people who would have used the metro switching to shared bikes, often using longer or unfamiliar routes when their regular stations were inaccessible~\cite{Yang2022}. Similarly, during a major flood in Porto Alegre, Brazil, the local \ac{bss} continued to function and enabled many daily activities to proceed despite widespread transport disruption~\cite{Araujo2025}. Evidence from Washington D.C. further shows that bike-sharing usage rises significantly during planned Metro shutdowns, demonstrating its role as an alternative mode in complete transit outages~\cite{Cheng2022}. Last, the COVID-19 pandemic further reinforced this role of \acp{bss} as a transport mode that brings resilience to the transport network. While \ac{pt} usage remained far below pre-pandemic levels, bike-sharing volumes nearly returned to normal once restrictions eased, with some even increasing its trips numbers~\cite{Teixeira2024}. For instance, in Barcelona's \textit{Bicing}, trip levels exceeded pre-pandemic counts within just a couple of months of reopening after the quarantine~\cite{Grau-Escolano2024}.